\section[Интерполирование функций, полиномы Чебышева, интерполяция с кратными узлами, кубические сплайны]{Интерполирование функций, полиномы Чебышева,\\ интерполяция с кратными узлами, кубические сплайны}
\label{sec:interpolation-chebyshev-hermite-splines}

\subsection{Полиномиальная интерполяция}

Пусть заданы попарно различные узлы
\[
x_0,x_1,\ldots,x_n
\]
и значения $y_i=f(x_i)$. \textbf{Задача интерполяции} состоит в построении функции $P$, удовлетворяющей
\[
P(x_i)=y_i,\qquad i=0,\ldots,n.
\]

\textbf{Теорема о существовании и единственности.} Для любых $n+1$ попарно различных узлов и значений $y_i$ существует единственный алгебраический многочлен $P_n$ степени не выше $n$, интерполирующий эти данные.

Действительно, система для коэффициентов в степенном базисе имеет матрицу Вандермонда, а
\[
\det V=\prod_{0\le i<j\le n}(x_j-x_i)\ne0.
\]

\subsection{Формы Лагранжа и Ньютона}

\textbf{Формула Лагранжа}:
\[
P_n(x)=\sum_{i=0}^{n}y_iL_i(x),
\qquad
L_i(x)=\prod_{\substack{k=0\\k\ne i}}^{n}
\frac{x-x_k}{x_i-x_k},
\]
причём $L_i(x_j)=\delta_{ij}$.

Для последовательного добавления узлов удобна \textbf{форма Ньютона}. Разделённые разности определяются рекурсивно:
\[
f[x_i]=f(x_i),
\qquad
f[x_i,\ldots,x_{i+m}]
=
\frac{f[x_{i+1},\ldots,x_{i+m}]-f[x_i,\ldots,x_{i+m-1}]}
{x_{i+m}-x_i}.
\]
Тогда
\[
\begin{split}
P_n(x)=&\,f[x_0]+f[x_0,x_1](x-x_0)+\cdots\\
&+f[x_0,\ldots,x_n]\prod_{i=0}^{n-1}(x-x_i).
\end{split}
\]
Для попарно различных узлов разделённая разность симметрична по своим аргументам; в частности,
\[
f[x_0,\ldots,x_n]
=
\sum_{j=0}^{n}
\frac{f(x_j)}{\displaystyle\prod_{\substack{i=0\\i\ne j}}^{n}(x_j-x_i)}.
\]
На равномерной сетке разделённые разности выражаются через конечные разности.

\subsection{Остаточный член и выбор узлов}

Если $f\in C^{n+1}[a,b]$, все узлы и точка $x$ лежат в $[a,b]$, то существует $\xi=\xi(x)\in(a,b)$ такое, что
\[
\boxed{
f(x)-P_n(x)
=
\frac{f^{(n+1)}(\xi)}{(n+1)!}
\omega_{n+1}(x),
\qquad
\omega_{n+1}(x)=\prod_{i=0}^{n}(x-x_i).
}
\]
Поэтому выбор узлов связан с уменьшением $\max|\omega_{n+1}|$.

На $[-1,1]$ \textbf{полиномы Чебышева}
\[
T_m(x)=\cos(m\arccos x)
\]
удовлетворяют $|T_m(x)|\le1$. Их нули
\[
x_k=\cos\frac{(2k+1)\pi}{2m},
\qquad k=0,\ldots,m-1,
\]
используются как интерполяционные узлы. Среди приведённых многочленов степени $m$ полином
\[
2^{1-m}T_m(x)
\]
имеет наименьшую равномерную норму на $[-1,1]$. Поэтому при $m=n+1$ чебышевские узлы минимизируют максимальную величину корневого многочлена и существенно уменьшают краевые осцилляции по сравнению с равномерной сеткой.

Следует отличать теорему Вейерштрасса о существовании хороших полиномиальных приближений от сходимости конкретного интерполяционного процесса. \textbf{Теорема Фабера} показывает, что не существует фиксированной стратегии выбора узлов, для которой интерполяционные полиномы сходились бы равномерно к каждой функции из $C[a,b]$. Для достаточно гладких функций и удачно выбранных узлов, в частности чебышевских, интерполяция сходится.

\subsection{Интерполяция с кратными узлами. Полином Эрмита}

Если кроме значений функции заданы значения производных, возникает \textbf{интерполяция Эрмита}. Пусть в $n+1$ различных узлах известны
\[
f(x_i)=y_i,
\qquad
f'(x_i)=y_i',
\qquad i=0,\ldots,n.
\]
Существует единственный многочлен $H_{2n+1}$ степени не выше $2n+1$, удовлетворяющий всем $2n+2$ условиям.

Пусть $L_i$ --- базис Лагранжа. Тогда удобная запись:
\[
H_{2n+1}(x)
=
\sum_{i=0}^{n}
\left[y_i h_i(x)+y_i'\widehat h_i(x)\right],
\]
где
\[
h_i(x)=
\left[1-2L_i'(x_i)(x-x_i)\right]L_i^2(x),
\qquad
\widehat h_i(x)=(x-x_i)L_i^2(x).
\]
Эти функции удовлетворяют
\[
h_i(x_j)=\delta_{ij},\quad h_i'(x_j)=0,
\qquad
\widehat h_i(x_j)=0,\quad \widehat h_i'(x_j)=\delta_{ij}.
\]

Если $f\in C^{2n+2}[a,b]$, то
\[
\boxed{
f(x)-H_{2n+1}(x)
=
\frac{f^{(2n+2)}(\xi)}{(2n+2)!}
\left(\prod_{i=0}^{n}(x-x_i)\right)^2.
}
\]

В общем случае, если в узле $x_i$ заданы производные до порядка $m_i-1$, полная кратность равна
\[
N=\sum_i m_i,
\]
и интерполяционный многочлен имеет степень не выше $N-1$.

\subsection{Кубические интерполяционные сплайны}

\textbf{Кубический интерполяционный сплайн} $S$ по узлам
\[
a=x_0<x_1<\cdots<x_n=b
\]
--- функция, для которой:
\begin{enumerate}
\item $S|_{[x_{i-1},x_i]}$ --- многочлен степени не выше $3$;
\item $S\in C^2[a,b]$;
\item $S(x_i)=y_i$ во всех узлах;
\item заданы ещё два краевых условия.
\end{enumerate}

Обозначим
\[
h_i=x_i-x_{i-1},
\qquad M_i=S''(x_i).
\]
На $[x_{i-1},x_i]$ сплайн имеет представление
\[
\begin{split}
S_i(x)=&\frac{M_{i-1}(x_i-x)^3}{6h_i}
+\frac{M_i(x-x_{i-1})^3}{6h_i}\\
&+\left(y_{i-1}-\frac{M_{i-1}h_i^2}{6}\right)\frac{x_i-x}{h_i}
+\left(y_i-\frac{M_ih_i^2}{6}\right)\frac{x-x_{i-1}}{h_i}.
\end{split}
\]
Непрерывность $S'$ приводит к трёхдиагональной системе
\[
\boxed{
\frac{h_i}{6}M_{i-1}
+\frac{h_i+h_{i+1}}{3}M_i
+\frac{h_{i+1}}{6}M_{i+1}
=
\frac{y_{i+1}-y_i}{h_{i+1}}
-
\frac{y_i-y_{i-1}}{h_i},
}
\]
$i=1,\ldots,n-1$.

Наиболее употребительные варианты двух дополнительных краевых условий:
\begin{itemize}
\item \textbf{естественный сплайн}: $M_0=M_n=0$;
\item \textbf{закреплённый (clamped) сплайн}: задаются $S'(a)$ и $S'(b)$;
\item \textbf{not-a-knot}: обеспечивается непрерывность третьей производной в первых и последних внутренних узлах,
\[
S_1'''(x_1)=S_2'''(x_1),
\qquad
S_{n-1}'''(x_{n-1})=S_n'''(x_{n-1});
\]
\item для периодических данных ($y_0=y_n$): $S'(a)=S'(b)$ и $S''(a)=S''(b)$.
\end{itemize}

\textbf{Теорема.} Для любых значений $y_i$ в строго упорядоченных узлах существует единственный естественный кубический интерполяционный сплайн.

Кроме того, естественный сплайн характеризуется вариационным свойством: среди функций $v\in H^2(a,b)$, принимающих те же значения в узлах, он единственным образом минимизирует ``энергию кривизны''
\[
\boxed{
\int_a^b (S''(x))^2\,dx
\le
\int_a^b (v''(x))^2\,dx.
}
\]
Если $f\in C^4[a,b]$ и краевые условия согласованы с $f$, то на квазиравномерной сетке типично $\|f^{(k)}-S^{(k)}\|_\infty=O(h^{4-k})$ для $k=0,1,2,3$. Естественные условия $S''(a)=S''(b)=0$ могут быть несогласованы с произвольной $f$, поэтому точность производных у концов в таком случае может ухудшаться.

\subsection{Что важно сказать на экзамене}

Нужно знать четыре блока:
\[
\boxed{
\text{Лагранж/Ньютон + остаток}
\rightarrow
\text{узлы Чебышева}
\rightarrow
\text{Эрмит}
\rightarrow
\text{кубический сплайн}}
\]
и уметь сформулировать теорему единственности интерполяционного многочлена, остаточный член, смысл чебышевских узлов и систему для вторых производных кубического сплайна.

\newpage
